%!TEX root = ../template.tex
%%%%%%%%%%%%%%%%%%%%%%%%%%%%%%%%%%%%%%%%%%%%%%%%%%%%%%%%%%%%%%%%%%%%
%% chapter2.tex
%% NOVA thesis document file
%%
%% Chapter with the template manual
%%%%%%%%%%%%%%%%%%%%%%%%%%%%%%%%%%%%%%%%%%%%%%%%%%%%%%%%%%%%%%%%%%%%

\typeout{NT FILE chapter2.tex}%

\chapter{Fundamental Concepts}
\label{cha:concepts}

In this section we aim to introduce some of the concepts that are relevant to the work
presented in this thesis. These concepts include Machine Learning techniques,
mainly toward Supervised Learning, and some very basic fundamentals about image types
that are relevant to our thesis.

\section{Machine Learning}
\gls{ml} is a field of study in computer science that utilizes
statistical techniques to enable computer systems to learn from data, recognize
patterns, and make predictions on new data. \gls{ml} is broadly categorized into
supervised and unsupervised learning. The choice between these categories
depend on the problem at hand; supervised learning is used when labeled
data is available, while unsupervised learning is used for finding hidden
patterns in unlabeled data as seen in Figure \ref{fig:ml-tree}.

For our case, supervised learning is more suitable, as we aim to classify
vegetation risk based on labeled aerial photographs/satellite imagery
in accordance with \gls{adai}'s risk classification standards, which is considered
an image classification problem.

\begin{figure}[ht]
    \centering
    \includegraphics[width=0.8\textwidth]{../Figures/MLTaxonomy.drawio.png}
    \caption{Machine Learning taxonomy inspired by \cite{image_taxonomy}.}
    \label{fig:ml-tree}
\end{figure}

When discussing \gls{ml}, it is essential to address feature extraction, normalization,
and standardization.

Feature extraction involves identifying and extracting the most relevant
features from the data. Techniques such as Principal Component Analysis
(PCA) and \gls{cnn} are commonly used for
this purpose. In the context of satellite/aerial imagery, important features
might include vegetation indices, texture, and spectral bands. These
features help in distinguishing different vegetation types \cite{rs13081411} and therefore
the associated fire risks.

Normalization and standardization are crucial steps in \gls{ml}, particularly
when dealing with diverse data types and scales. They ensure that each
feature contributes equally to the model, improving the accuracy and
robustness of the classifier.
For satellite imagery, features can vary significantly in scale, which can pose challenges for
classification algorithms \cite{Adegun2023}. Without normalization or standardization,
certain features might disproportionately influence the classification
process, leading to worse results.

\subsection{Supervised Learning}

Supervised Learning (SL) models learn to map input data to the correct output
by minimizing the error between the predicted and actual values. The goal of
SL is to generalize from the training data to make accurate
predictions on new, unseen data \cite{geeksforgeeks_supervised_learning}.

There are two main types of supervised learning: regression and classification.
In regression, the goal is to predict a continuous value
based on input features. For example, predicting house prices based on features
such as location, size, and number of rooms.

In classification, the goal is to predict a discrete
class label based on input features. For example, our problem at hand, classifying
house fire risk based on the surrounding vegetation features, is a classification problem.

In supervised learning, the dataset is divided into two parts: the training set
and the test set. The training set is used to train the model, while the test
set is used to evaluate the model's performance on unseen data. The model is
trained by adjusting its parameters to minimize the error between the predicted
and actual values on the training set. Once the model is trained, it is evaluated
on the test set to assess its performance and generalization capabilities.

\subsubsection{Evaluation}

To evaluate the performance of a supervised learning model, various metrics are used,
such as accuracy, precision, recall, F1 score, and kappa statistic. These metrics help
in assessing the model's performance, identifying areas for improvement, and comparing
different models.

The accuracy of a model is the proportion of correct predictions to the total number
of predictions \eqref{eq:accuracy}. While simple and intuitive, accuracy can be misleading
when dealing with imbalanced datasets. Precision measures the proportion of true positive
predictions to the total number of positive predictions \eqref{eq:precision}, which is useful
when the cost of false positives is high. Recall, also known as sensitivity, measures the
proportion of true positive predictions to the total number of actual positive instances
\eqref{eq:recall}, making it useful when the cost of false negatives is high. The F1 score, the
harmonic mean of precision and recall \eqref{eq:f1}, balances the two metrics and provides a
more comprehensive performance measure \cite{vidalwadi2024}.


\begin{equation}
    Accuracy = \frac{TP + TN}{TP + TN + FP + FN}
    \label{eq:accuracy}
\end{equation}

\begin{equation}
    Precision = \frac{TP}{TP + FP}
    \label{eq:precision}
\end{equation}

\begin{equation}
    Recall = \frac{TP}{TP + FN}
    \label{eq:recall}
\end{equation}

\begin{equation}
    F1 = 2 \times \frac{Precision \times Recall}{Precision + Recall}
    \label{eq:f1}
\end{equation}
In these equations:
\begin{itemize}
    \item True Positives (TP): The number of correctly predicted positive instances.
    \item True Negatives (TN): The number of correctly predicted negative instances.
    \item False Positives (FP): The number of instances incorrectly predicted as positive.
    \item False Negatives (FN): The number of instances incorrectly predicted as negative.
\end{itemize}

In addition to these metrics, the kappa statistic \eqref{eq:kappa} is used to evaluate the agreement
between predicted and actual classifications, accounting for the possibility of agreement
occurring by chance. It is especially useful for imbalanced datasets. Kappa values range
from -1 to 1, where 1 indicates perfect agreement, 0 indicates no agreement beyond chance,
and negative values indicate disagreement \cite{kappa}. The kappa statistic is calculated as follows:
\begin{equation}
    \kappa = \frac{P_o - P_e}{1 - P_e}
    \label{eq:kappa}
\end{equation}
where \(P_o\) is the observed agreement and \(P_e\) is the expected agreement by chance.

Including the kappa statistic in the evaluation framework provides a more reliable assessment
of model performance, ensuring that the model's effectiveness is thoroughly evaluated.

For multiclass classification problems, such as our own, evaluation metrics can be extended
to provide a more comprehensive assessment. One approach is macro-averaging, where the
metric is calculated independently for each class and then averaged. This method treats
all classes equally. Another approach is micro-averaging, which aggregates the contributions
of all classes to compute the average metric, treating all instances equally. Additionally,
weighted-averaging calculates the metric independently for each class and then takes the
average, weighted by the number of true instances for each class. These methods ensure that
the evaluation metrics accurately reflect the performance across all classes.


\subsection{Unsupervised Learning}

Unsupervised Learning (UL) main goal is to discover hidden patterns
and structures within the data without explicit guidance.

Unsupervised learning is particularly useful for tasks such as clustering, dimensionality
reduction, and anomaly detection. For our problem, unsupervised learning does not
apply directly, as we have labeled data for the fire risk classification.

\section{Deep Learning on Artificial Neural Networks}

\gls{dl} is a specialized subset of machine learning that makes use of
artificial neural networks (ANNs) to identify intricate patterns within large datasets.
Unlike traditional machine learning algorithms, DL algorithms can learn directly
from raw data, automatically extract features, and make highly accurate predictions.
This capability has significantly contributed to various fields, including computer vision,
natural language processing, and speech recognition.

DL models consist of multiple layers of artificial neurons that process
input data and learn hierarchical representations of this data \cite{spotfire_neural_network}. The most prevalent
type of \gls{dl} model is the deep neural network, which is structured with an
input layer, one or more hidden layers, and an output layer as seen in Figure \ref{fig:dnn-example}. Each layer contains
numerous neurons that perform computations on the input data and pass the results
to the subsequent layer. The hidden layers are responsible for learning progressively
complex features of the data, whereas the output layer generates the final prediction.

\begin{figure}[ht]
    \centering
    \includegraphics[width=0.8\textwidth]{../Figures/DNN.png}
    \caption{Artificial Neural Network Example from \cite{deep_neural_network_img}.}
    \label{fig:dnn-example}
\end{figure}

The flexibility and power of deep neural networks make them highly suitable for a
variety of complex tasks. Moreover, when it comes to image data and image classification,
a specific type of \gls{dl} model, known as \gls{cnn},
has proven to be exceptionally effective \cite{rs13224712}.

\subsection{Multilayer Perceptrons}

\gls{mlp} are one of the most basic forms of artificial neural
networks. They consist of multiple layers of neurons, including an input layer, one
or more hidden layers, and an output layer. \gls{mlp} are fully connected, meaning each
neuron in one layer connects to every neuron in the next layer as seen in Figure
\ref{fig:dnn-example}.

\subsubsection{Activation Functions}

A critical component of \gls{mlp} is the activation function used in the neurons. Activation
functions introduce non-linearity into the network, enabling it to learn complex patterns.
Nowadays the most popular choice for activation functions is the Rectified Linear Unit (ReLU) \ref{eq:relu},
due to its simplicity and good performance. Other common activation functions include the
Sigmoid \ref{eq:sigmoid}, and Tanh (Hyperbolic Tangent) \ref{eq:tanh}.

\begin{equation}
    f(x) = \max(0, x)
    \label{eq:relu}
\end{equation}

\begin{equation}
    f(x) = \frac{1}{1 + e^{-x}}
    \label{eq:sigmoid}
\end{equation}

\begin{equation}
    f(x) = \tanh(x) = \frac{e^x - e^{-x}}{e^x + e^{-x}}
    \label{eq:tanh}
\end{equation}
\smallskip

While \gls{mlp} are powerful for certain types of data, they are not as efficient as \gls{cnn} for
image classifications due to their inability to exploit the spatial structure of images \cite{zhang2021dive},
at least, not on their own.


\subsubsection{Forward Propagation}

Forward propagation is the process by which input data passes through the layers of the
MLP to generate an output. In each layer, the input is transformed linearly using weights
and biases, followed by the application of an activation function.

\subsubsection{Backpropagation}

Backpropagation is the basic learning algorithm used to train \gls{mlp}. It involves two main
steps: the forward pass and the backward pass. In the backward pass, the network
calculates the gradient of the loss function with respect to each weight by the
chain rule, propagating errors backward from the output layer to the input layer.
This gradient is then used to update the weights to minimize the loss function,
typically through gradient descent optimization \cite{zhang2021dive}.

\subsubsection{Dropout}

Dropout is a regularization technique used in fully connected layers to prevent
overfitting. In standard dropout regularization, a fraction of the nodes in each
layer is zeroed out, and each layer is then debiased by normalizing by the fraction
of nodes that were retained as seen in Figure \ref{fig:dropout-example}. Specifically, with dropout
probability \( p \), each intermediate activation \( h \) is replaced by a random
variable \( h' \) as follows:

\[
    h' =
    \begin{cases}
        0             & \text{with probability } p   \\
        \frac{h}{1-p} & \text{with probability } 1-p
    \end{cases}
\]

During training, this process prevents the network from becoming too reliant on
specific neurons, encouraging it to learn more robust features. Typically, dropout
is disabled at test time. When evaluating a trained model with a new example, we
do not drop out any nodes and thus do not need to normalize.

\begin{figure}[ht]
    \centering
    \includegraphics[width=0.8\textwidth]{../Figures/dropout-ex.png}
    \caption{MLP dropout example taken from \cite{zhang2021dive}.}
    \label{fig:dropout-example}
\end{figure}

However, there are some exceptions to this rule. Some researchers \cite{zhang2021dive} use dropout at
test time as a heuristic for estimating the uncertainty of neural network
predictions. By running the network multiple times with dropout enabled and
observing the variance in predictions, one can gauge the confidence of the
network. If the predictions agree across many different dropout outputs, then
the network is considered to be more confident in its predictions.


\subsection{Convolutional Neural Networks}

\gls{cnn} are a type of deep neural network that
is particularly well-suited for image classification tasks. \cite{rs-survey-cnns} \gls{cnn} are designed
to automatically learn spatial hierarchies of features from images, making them
ideal for tasks such as object detection, image segmentation, and facial recognition.

In the context of the project, since our goal is to classify the fire risk of isolated houses
based on parameters like the proximity and type of surrouding vegetation. This requires
identifying specific objects such as the house and the distance between the house and
the specific types of vegetation. \gls{cnn} are well-suited for this task due to their ability
to automatically learn and detect spatial features from images.

\gls{cnn} utilize three main types of layers to process image data, namely Convolutional layers,
Pooling layers, and Fully Connected layers. The following figure is
an example of a simple CNN architecture \ref{fig:cnn-architecture}.

\begin{figure}[ht]
    \centering
    \includegraphics[width=0.8\textwidth]{../Figures/example-cnn.png}
    \caption{Architecture of \gls{cnn} from \cite{cheng2020remote}.}
    \label{fig:cnn-architecture}
\end{figure}


\subsubsection{Convolutional Layers}
Convolutional layers are the core building blocks of \gls{cnn}, responsible for feature
extraction from input images. These layers apply learnable filters to the input image,
which slide across the image's width and height to produce feature maps. Each feature
map highlights specific features such as edges, textures, or patterns. By stacking
multiple convolutional layers, the network can learn increasingly complex feature
representations, from simple edges in the initial layers to intricate shapes and
objects in the deeper layers.

The input to the convolutional layers \( X \in \mathbb{R}^{n \times w \times h} \) consists
of \( n \) 2-D feature maps of size \( w \times h \). The output
\( H \in \mathbb{R}^{m \times w' \times h'} \) of the convolutional layers is \( m \) 2-D
feature maps of size \( w' \times h' \) produced via convolutional kernels \( W \). Here,
\( W \in \mathbb{R}^{m \times l \times l \times n} \) represents \( m \) trainable filters
of size \( l \times l \times n \) (typically \( l = 1, 3, \) or \( 5 \)). The entire process
of convolution is described as follows:

\begin{equation}
    H = f(W * X + b)
    \label{eq:convolution}
\end{equation}
where \( * \) denotes the 2-D convolution operation, and \( b \) represents the
\( m \)-dimensional bias term. In general, a nonlinear activation function
\( f \) is applied after the convolution operation. As the convolutional structure
deepens, the convolutional layers can capture different levels of features from the
input feature maps.


\subsubsection{Pooling Layers}
Pooling layers are used to reduce the spatial dimensions of the feature maps, which
helps to decrease the computational load and the number of parameters in the network \cite{poolinglayers-ref}.
The two most common types of pooling are max pooling, which
partitions the input feature map into a set of non-overlapping regions and
outputs the maximum value from each region, and average pooling, which
outputs the average value within each region \ref{fig:max-pool-example}.

\begin{figure}[ht]
    \centering
    \includegraphics[width=1\textwidth]{../Figures/pooling-exs.png}
    \caption{Max Pooling and Average Pooling examples from \cite{vedantkumar}.}
    \label{fig:max-pool-example}
\end{figure}

The pooling operation provides several benefits:
\begin{itemize}
    \item \textbf{Dimensionality Reduction:} By reducing the size of the feature maps, pooling layers help to decrease the computational cost and the risk of overfitting.
    \item \textbf{Translation Invariance:} Makes the network more robust to small translations or distortions in the input image. This means that the network can better recognize objects regardless of their location within the image.
\end{itemize}

\subsubsection{Fully Connected Layers}
Fully connected layers, also known as dense layers, are used towards the end of a
CNN to integrate the high-level features extracted by the convolutional and pooling
layers. In these layers, each neuron is connected to every neuron in the previous
layer like an MLP, allowing for the combination and processing of features to produce the final
output. This can be a class score for classification tasks, bounding box coordinates
for object detection, or pixel-wise labels for segmentation tasks. Fully connected
layers typically use activation functions like ReLU to introduce non-linearity, enabling
the network to learn complex mappings between inputs and outputs.

These layers process their input \( \tilde{X} \) with
a linear transformation by weights \( \tilde{W} \) and bias \( \tilde{b} \), and then
map the output of the linear transformation through a nonlinear activation function \( f \).
The entire process can be formulated as:

\begin{equation}
    y = f(\tilde{W} \cdot \tilde{X} + \tilde{b})
    \label{eq:fully-connected}
\end{equation}

In classification tasks, to output the probability of each class, a softmax classifier
is typically connected to the last fully connected layer. The softmax classifier
normalizes the output \( y \in \mathbb{R}^c \) (where \( c \) is the number of classes)
to a probability distribution between 0 and 1. This can be described as:

\begin{equation}
    P(y_i) = \frac{e^{y_i}}{\sum_{j=1}^{c} e^{y_j}}
    \label{eq:softmax}
\end{equation}

The output of the softmax classifier represents the probability that the given input
image belongs to each class. To avoid overfitting, the dropout method can be applied
to the fully connected layers, as these layers usually contain a large number of
parameters.


\subsubsection{Modern CNN Architectures}
Modern \gls{cnn} have significantly advanced the field
of computer vision and are highly suitable for complex image analysis tasks such as our
classification problem at hand. By leveraging the unique capabilities of existing advanced
architectures, we can help contributing for our fire risk assessment model.

Some of the most popular CNN architectures include AlexNet \cite{krizhevsky2012imagenet}, VGGNet
\cite{simonyan2014very}, ResNet \cite{he2016deep}, DenseNet \cite{huang2017densely} and Inception \cite{szegedy2015going},
they have introduced several innovations that can be beneficial for our classification.
For example, ResNet learns residual functions with reference to the layer inputs rather
than learning unreferenced functions. Instead of hoping each few stacked layers directly
fit a desired underlying mapping, residual networks let these layers fit a residual mapping.
They stack residual blocks to form a deep network that can be trained using standard
backpropagation techniques. This architecture has proven to be highly effective in training
very deep neural networks \cite{DBLP:journals/corr/HeZRS15}.

While image segmentation could potentially aid in easier identification of critical
features like vegetation and structures around houses, it requires extensive labeled
data for training. Currently, our dataset is specifically labeled for fire risk classification
rather than detailed image segmentation. This limitation makes pursuing segmentation techniques
challenging without additional labeled data.

Given these considerations, our focus remains on leveraging the strengths of CNN architectures
for classification tasks related to fire risk assessment. This approach allows us to effectively
utilize the available labeled data to predict fire risk based on identified features without
the need for extensive segmentation.

\subsubsection{Transfer Learning}

Transfer learning is a machine learning technique where a model trained on one task and
is leveraged as a starting point for a model on a second related task. The idea behind transfer
learning is to transfer knowledge learned from a source domain (task) to a target domain (task)
to improve learning efficiency and performance, especially when the target dataset is small or
labeled data is limited. There are two main approaches to transfer learning:

\textbf{Feature Extraction:} In feature extraction, the learned features (representations)
from the pre-trained model are used directly as inputs for the new model. The early layers of
the pre-trained model, which capture generic features applicable to many tasks (like edges, textures),
are frozen (not trained) while the later layers are replaced or fine-tuned for the new task.

\textbf{Fine-tuning:} Fine-tuning extends feature extraction by unfreezing and jointly training
some or all layers of the pre-trained model along with the new task-specific layers. This allows
the model to adapt more specifically to the nuances of the target task while still benefiting from
the pre-trained weights and learned features.

Transfer learning is widely used in \gls{dl}, particularly with \gls{cnn}
trained on large-scale datasets like ImageNet \cite{deng2009imagenet}. By initializing a model with
weights pre-trained on such datasets, transfer learning accelerates convergence and improves
generalization on new datasets with fewer labeled examples. This approach has been pivotal in
achieving state-of-the-art results in various domains, including image classification, object
detection, and natural language processing.

\subsection{Ensemble Methods}

In addition to the modern CNN architectures discussed, ensemble methods can provide
significant improvements in classification tasks, particularly in complex domains
such as our fire risk assessment. Ensemble methods work by combining multiple models
to produce a single, more accurate model. The idea is that by aggregating the
predictions of multiple models, the overall performance can be enhanced, leading to
better generalization and robustness.

There are several ways ensembles can be applied. One example, that could be relevant to our problem is training
models separately on different parameters. For instance, in fire risk classification,
which is influenced by various factors such as the distance of a house from different
types of vegetation, training individual models for each parameter can improve the
overall classification accuracy. This approach helps pinpoint issues in the classification
process and enhances the overall model's performance.

Using a single model to handle all the specific parameters might struggle to
find patterns due to the complexity of the data. By training separate CNN models for
different identification tasks and then combining their outputs in a final MLP
(Multi-Layer Perceptron) model, we can more effectively classify the fire risk of a house.
This ensemble approach allows each CNN to specialize in a specific aspect of the data, and
the final MLP can integrate these specialized insights to produce a more accurate and
reliable classification.



\subsection{Autoencoders}

Autoencoders are a type of artificial neural network used for unsupervised learning
of efficient codings. The primary goal of an autoencoder is to learn a representation
(encoding) for a set of data, typically for the purpose of dimensionality reduction. An
autoencoder consists of two main components: the encoder and the decoder \cite{goodfellow2016deep}.

The encoder compresses the input into a latent-space representation, which is a
lower-dimensional space that captures the essential features of the input data.
This latent space can be thought of as a bottleneck where only the most important
aspects of the data are retained. The process of encoding can be formulated as:
\begin{equation}
    h = f (W \cdot x + b)
    \label{eq:enconding}
\end{equation}
where \(h \in \mathbb{R}^n\) is the output of the hidden layers, \(f\) denotes a
nonlinear mapping, \(W \in \mathbb{R}^{n \times m}\) stands for the encoding weight
matrix, \(x \in \mathbb{R}^m\) denotes the input of the autoencoder, and \(b \in \mathbb{R}^n\)
is the bias vector \cite{cheng2020remote}.

The decoder then attempts to reconstruct the original input data from this compressed
representation. The decoding process is the inverse of encoding and can be formulated as:
\begin{equation}
    \tilde{x} = f (W' \cdot h + b')
    \label{eq:decoding}
\end{equation}
where \(\tilde{x} \in \mathbb{R}^m\) represents the reconstructed output,
\(W' \in \mathbb{R}^{m \times n}\) is the decoding weight matrix, and \(b' \in \mathbb{R}^m\)
stands for the bias vector. The network is trained to minimize the difference between the
input data and the reconstructed output, often using a loss function such as mean squared
error. Through this process, the autoencoder learns to extract the most significant
features of the input data.

Autoencoders have wide-ranging applications, including anomaly detection, data compression,
image processing, and feature extraction. They are a powerful tool enabling the discovery
of meaningful patterns and structures within data without the need for labeled examples.

\subsection{Generative Adversarial Networks}

\gls{gan} are a powerful and innovative approach in the
field of machine learning, particularly for generating realistic data. \gls{gan} consist of
two neural networks, the generator (G) and the discriminator (D), which are trained
simultaneously through a process known as adversarial training as seen in Figure \ref{fig:gan}.

The generator's role is to create synthetic data that resembles real data, while the
discriminator's job is to distinguish between real and synthetic data.
Over time, both networks improve in their respective tasks, leading to the generation
of highly realistic data by the generator.

\begin{figure}[h]
    \centering
    \includegraphics[width=0.8\textwidth]{../Figures/GAN-ARCH.png}
    \caption{Illustration of GAN architecture with Generator (G) and Discriminator (D) from \cite{app12147075}.}
    \label{fig:gan}
\end{figure}

The generator takes a random noise vector \(z\) from a prior distribution \(p_z(z)\) and
transforms it into data samples \(G(z; \theta_g)\), where \(\theta_g\) represents the
parameters of the generator. The discriminator, on the other hand, evaluates these
generated samples along with real samples from the true data distribution \(p_{data}(x)\) and
outputs a probability \(D(x; \theta_d)\) that the input is a real data sample, with
\(\theta_d\) representing the parameters of the discriminator.

The training of \gls{gan} involves a minimax game with the following objective function:

\begin{equation}
    \min_G \max_D V(G, D) = \mathbb{E}_{x \sim p_{data}(x)}[\log D(x)] + \mathbb{E}_{z \sim p_z(z)}[\log(1 - D(G(z)))]
\end{equation}

Here, the discriminator \(D\) aims to maximize the probability of correctly identifying
real and fake samples, while the generator \(G\) aims to minimize the discriminator's
ability to distinguish between the two \cite{NIPS2014_5ca3e9b1}. This adversarial process continues until the
discriminator cannot differentiate between real and synthetic data with better than random
guessing accuracy.

\gls{gan} have been successfully applied in various fields, such as image generation,
style transfer, and data augmentation. Despite their success, training them can be
challenging due to issues like mode collapse and instability.

They have opened new possibilities for
generating high-quality synthetic data that closely mimics real-world data, making them
a significant tool in the arsenal of modern machine learning methods.


\subsection{Loss Functions}
Loss functions measure the difference between predicted and actual values during model training.
These functions play a crucial role in guiding the optimization
process towards improving model accuracy and generalization. The choice of an appropriate
loss function depends on the nature of the task, whether it involves regression, binary
classification, or multi-class classification. This section explores several key loss
functions commonly used in the training of neural networks.

\subsubsection{Mean Squared Error (MSE)}
Mean Squared Error is widely used for regression tasks. It calculates the average
of the squared differences between predicted and actual values:
\begin{equation}
    \text{MSE} = \frac{1}{n} \sum_{i=1}^{n} (y_i - \hat{y}_i)^2
    \label{eq:mse}
\end{equation}
where \( y_i \) is the actual value and \( \hat{y}_i \) is the predicted value. MSE
is sensitive to outliers due to the squaring term, making large errors more significant.

\subsubsection{Cross-Entropy Loss}
Cross-Entropy Loss, also known as log loss, is used for classification tasks.
It measures the performance of a classification model whose output is a probability
value between 0 and 1:

\begin{equation}
    \text{Cross-Entropy Loss} = -\frac{1}{n} \sum_{i=1}^{n} [y_i \log(\hat{y}_i) + (1 - y_i) \log(1 - \hat{y}_i)]
    \label{eq:cross-entropy}
\end{equation}
where \( y_i \) is the actual binary label and \( \hat{y}_i \) is the predicted probability.
For multi-class classification, a generalized version called categorical cross-entropy is used.

\subsubsection{Categorical Cross-Entropy Loss}
Categorical Cross-Entropy Loss is used for multi-class classification tasks where each
sample can belong to one of several classes, aligning to our problem.
It is an extension of the binary cross-entropy loss:

\begin{equation}
    \text{Categorical Cross-Entropy Loss} = -\sum_{i=1}^{n} \sum_{j=1}^{k} y_{ij} \log(\hat{y}_{ij})
    \label{eq:categorical-cross-entropy}
\end{equation}
where \( y_{ij} \) is a binary indicator (0 or 1) if class label \( j \) is the correct classification for observation \( i \), \( \hat{y}_{ij} \) is the predicted probability that observation \( i \) belongs to class \( j \), and \( k \) is the number of classes. This loss function encourages the predicted probability to be high for the correct class and low for all other classes.


\subsection{Optimization Algorithms}
Optimization algorithms play a crucial role in training \gls{dl} models by
minimizing the loss function and adjusting the model's parameters. These algorithms
iteratively update the weights of the neural network to find the set of parameters
that best fits the training data. The choice of optimization algorithm can significantly
impact the model's performance and training efficiency \cite{ruder2016overview}.

\subsubsection{Gradient Descent}

Gradient Descent is one of the most fundamental optimization algorithms. It works
by calculating the gradient of the loss function with respect to the model parameters
and then updating the parameters in the opposite direction of the gradient. The update
rule for Gradient Descent is given by:

\begin{equation}
    \theta = \theta - \eta \nabla_\theta J(\theta)
\end{equation}
where \( \theta \) represents the model parameters, \( \eta \) is the learning rate,
and \( \nabla_\theta J(\theta) \) is the gradient of the loss function \( J \) with respect
to \( \theta \). While effective, standard Gradient Descent can be slow and may get stuck
in local minima.

\subsubsection{Stochastic Gradient Descent (SGD)}

Stochastic Gradient Descent (SGD) is a variation of Gradient Descent that updates the
model parameters using a single training example or a small batch of examples at each
iteration. This introduces noise into the optimization process, which can help the
algorithm escape local minima and converge faster. The update rule for SGD is:

\begin{equation}
    \theta = \theta - \eta \nabla_\theta J(\theta; x^{(i)}, y^{(i)})
\end{equation}
where \( (x^{(i)}, y^{(i)}) \) represents a single training example. While SGD can converge
faster than batch Gradient Descent, it may result in a noisier optimization process.

\subsubsection{Adam}

The Adam (Adaptive Moment Estimation) optimizer combines the advantages of two other
popular optimization algorithms: AdaGrad and RMSProp \cite{ruder2016overview}. Adam maintains adaptive learning
rates for each parameter by computing the first and second moments of the gradients.
The update rule for Adam is:

\begin{equation}
    m_t = \beta_1 m_{t-1} + (1 - \beta_1) g_t
\end{equation}
\begin{equation}
    v_t = \beta_2 v_{t-1} + (1 - \beta_2) g_t^2
\end{equation}
\begin{equation}
    \hat{m}_t = \frac{m_t}{1 - \beta_1^t}
\end{equation}
\begin{equation}
    \hat{v}_t = \frac{v_t}{1 - \beta_2^t}
\end{equation}
\begin{equation}
    \theta_t = \theta_{t-1} - \eta \frac{\hat{m}_t}{\sqrt{\hat{v}_t} + \epsilon}
\end{equation}
where \( g_t \) is the gradient at time step \( t \), \( m_t \) and \( v_t \) are the first
and second moment estimates, \( \beta_1 \) and \( \beta_2 \) are the exponential decay rates
for these estimates, and \( \epsilon \) is a small constant for numerical stability. Adam is
widely used because of its efficiency and effectiveness in handling sparse gradients and
non-stationary objectives.


\section{Fundamentals of Aerial Photographs: RGB and IRG}
There are several types of images used in remote sensing, each capturing different
aspects of the electromagnetic spectrum. The two important types of images available for our
classification are \gls{rgb} and \gls{irg} images obtained from aerial flights and satellite imagery.

\subsection{RGB Images}
\gls{rgb} imagery is a standard image format that captures visible light using three color channels:
red, green, and blue. These images are similar to what the human eye perceives and are widely used
in various applications, including mapping, land cover classification, and environmental monitoring.
\gls{rgb} images provide valuable information about the visual characteristics of vegetation, such
as color and texture.

\subsection{IRG Images}
\gls{irg} images are a type of false-color composite that combines the green,
red, and near-infrared (NIR) bands of the electromagnetic spectrum \ref{fig:spectrum-example}. In \gls{irg} images, the NIR
band is assigned to the red channel, the green band to the green channel, and the red band
to the blue channel. This false-color combination is particularly advantageous for vegetation
analysis because plants reflect a significant amount of infrared light, which is not visible
to the human eye but can be detected by sensors. The infrared reflectance is a key indicator
of plant health and vigor, as healthy vegetation reflects more infrared light compared to
stressed or unhealthy plants.

\begin{figure}[ht]
    \centering
    \includegraphics[width=0.8\textwidth]{../Figures/NIR.png}
    \caption{Electromagnetic spectrum from \cite{nir_spectroscopy_nutshell}.}
    \label{fig:spectrum-example}
\end{figure}

The benefits of the \gls{irg} image for vegetation analysis are due to several technical reasons:

\textbf{Chlorophyll Reflectance Properties}:
Chlorophyll in plant leaves strongly absorbs light in the red and blue portions of the spectrum for use in
photosynthesis, but reflects green light, which is why plants appear green to our eyes. In the near-infrared
region, however, chlorophyll and the internal structure of the leaves cause significant reflection. This
reflectance increases with the density and health of the vegetation. By utilizing the NIR band, \gls{irg} images
can capture this high reflectance, highlighting areas of healthy vegetation more clearly than \gls{rgb} images as
seen in Figure \ref{fig:rgbxnirxnvdi}.

\textbf{Vegetation Indices}:
The combination of NIR and red bands is foundational for calculating vegetation indices,
such as the Normalized Difference Vegetation Index (NDVI) which can be formulated as \ref{eq:ndvi}.

\begin{equation}
    \text{NDVI} = \frac{\text{NIR} - \text{Red}}{\text{NIR} + \text{Red}}
    \label{eq:ndvi}
\end{equation}
where NIR and Red represent the reflectance values in the near-infrared and red bands, respectively.
NDVI has been widely used in remote sensing \cite{huang2021commentary} and it might prove
worthy to use for our classification problem, if we are having problems in differentiating the health
of the vegetation, as it provides a quantifiable measure of vegetation health. High NDVI values indicate healthy,
photosynthetically active vegetation, while lower values can indicate stressed or sparse vegetation. The \gls{irg} composite,
by emphasizing the NIR band, directly supports the visualization and calculation of such indices.

\begin{figure}[ht]
    \centering
    \includegraphics[width=1\textwidth]{../Figures/RGBxNIRxNVDI.png}
    \caption{\gls{rgb}, NIR and NDVI images of same crop field captured with rotary platform small unmanned aerial system (sUAS) from \cite{iastate2016imagery}.}
    \label{fig:rgbxnirxnvdi}
\end{figure}

\textbf{Stress Detection}:
Plants under stress (due to disease, drought, or other factors) exhibit reduced chlorophyll activity
and altered cell structure, leading to changes in their reflectance properties. Stressed plants typically
reflect less NIR light and more red light than healthy plants \cite{BUITRAGO201622}. The \gls{irg} images accentuate these differences,
allowing for the early detection of plant stress before visible symptoms appear, thus providing critical
information for vegetation monitoring and management.

Therefore, making use of \gls{irg} images can provide valuable insights to our classifier as
it can help in distinguishing between healthy and unhealthy vegetation, which is a critical
factor in the assessment of fire risk.